\documentclass[12pt]{article}
\usepackage[utf8]{inputenc}
\usepackage[T1]{fontenc}
\usepackage{hyperref}
\title{Wolof Language Resources for K-12 Curriculum Development}
\author{Compiled by Claude Sonnet 5.5 (claude-sonnet-5-5) for Andr\'as Kornai}
\date{}
\begin{document}
\maketitle

\section{Introduction}
Wolof (ISO 639-3 \texttt{wol}) is a language of wider communication, spoken mainly in Senegal, and also in the Gambia and Mauritania. Unlike many languages in this survey, it has dictionaries, grammars, university courses, Scripture and a national bilingual education policy in Senegal. The Gambian variety differs from the Senegalese one.

\section{Primary Educational Resources}
\begin{itemize}
\item The Senegalese government adopted a national bilingual education reform, MOHEBS (Mod\`ele harmonis\'e de l'enseignement bilingue au S\'en\'egal), in April 2021. According to the KIX/AIR brief, it is deployed from 2023--2024 in nine regions with six national languages. In the first two years about 80\% of time is in the national language. The brief and a related news item report that funding was weakened by the 2025 USAID withdrawal. \url{https://air.org/sites/default/files/2024-05/KIX-Strengthening-Bilingual-Education-Senegal-Brief-May-2024.pdf}
\item Columbia Teachers College blog on bilingual education research in Senegal: \url{https://www.tc.columbia.edu/cae/our-stories/blogs/bilingual-education-research-in-senegal/}
\item University of Florida Wolof course syllabi list \emph{Nanu Jang Wolof} (Mamarame Seck) and \emph{Nanu D\'egg Wolof} (Omar Ka) as course texts: \url{https://languages.ufl.edu/wp-content/uploads/sites/238/Spr21_WOL2201_Syllabus.pdf}
\end{itemize}

\section{Linguistic and Reference Materials}
\begin{itemize}
\item Sana Camara, \emph{Wolof Lexicon and Grammar} (2006), a Wolof--English and English--Wolof lexicon with grammar: \url{https://www.langcen.cam.ac.uk/resources/langw/data/wo_112.html}
\item Nyima Kantorek, \emph{Wolof--English/English--Wolof Dictionary \& Phrasebook}, based mainly on Gambian Wolof: \url{https://backstory.london/products/9780781810869}
\item Janga Wolof multilingual dictionaries: \url{https://payhip.com/PolyglotBooks/blog/articles/2-new-wolof-dictionaries-now-available}
\item Lexilogos Wolof dictionary page, which lists older and modern works such as Kob\`es 1923 and Munro and Gaye, \emph{Ay Baati Wolof}: \url{https://lexilogos.com/english/wolof_dictionary.htm}
\end{itemize}

\section{Religious and Community Resources}
\begin{itemize}
\item Complete Wolof Bible, \emph{K\`addug Y\`alla gi}, with online audio: \url{https://dbs.org/bibles/WOLKYG}
\item Scripture Earth index for Gambian Wolof (\texttt{wof}), which lists Gospels from 1907: \url{https://www.scriptureearth.org/00i-Scripture_Index.php?iso=wof}
\end{itemize}

\section{Digital and Audio Resources}
\begin{itemize}
\item Online audio for Gaye's \emph{Practical Course in Wolof}: \url{https://www.langmedia.fivecolleges.edu/online-audio-practical-course-wolof-gaye}
\item Global Recordings Network Wolof audio (evangelism and Bible teaching): \url{https://globalrecordings.net/en/language/wol}
\item iOS app ``Wolof Fran Engl Arabe'': \url{https://apps.apple.com/us/app/id6463724583}
\end{itemize}

\section{Curriculum Development Context}
Wolof is one of the Senegalese national languages used in the bilingual reform, so a curriculum could build on existing practice. ARED ran a simultaneous Wolof/Pulaar--French biliteracy programme from 2009 to 2018. Learner materials exist mostly for adult and university learners, and most of the dictionaries found are not aimed at schools. The current status of the reform's funding and national rollout was not verified. Gambian and Senegalese varieties differ, and materials should be matched to the variety.
\end{document}